\chapter{Comparación de objetivos planificados vs.\ realizados}
\label{apx:objetivos}

El objetivo general del trabajo fue diseñar e implementar un sistema multiagente
inteligente que integrara análisis predictivo, procesamiento de noticias y generación de
alertas personalizadas para el seguimiento de activos financieros.
La tabla~\ref{tab:apx-objetivos} compara lo planificado con lo efectivamente implementado.

\setlength{\extrarowheight}{6pt}
\begin{table}[H]
\centering
\caption{Comparación de objetivos planificados vs.\ realizados.}
\label{tab:apx-objetivos}
\small
\begin{tabular}{p{4.5cm}p{4.5cm}l}
\toprule
\textbf{Planificado} & \textbf{Realizado} & \textbf{Estado} \\
\midrule
5 agentes especializados & MarketAgent, ModelAgent, SentimentAgent, RecommendationAgent,
AlertAgent & Cumplido \\
\midrule
ML para series temporales & Ensemble de clasificación binaria (RF, GB, XGB, LightGBM):
Accuracy 57.0\%, F1 58.9\% & Cumplido \\
\midrule
NLP para análisis de sentimiento & Ensemble FinBERT (40\%) + VADER (25\%) + Lexicón (20\%) +
TextBlob (15\%) & Cumplido \\
\midrule
Interfaz gráfica web & Dashboard Streamlit con visualización de métricas y alertas &
Cumplido \\
\midrule
Evaluación con métricas de desempeño & 30 pruebas funcionales (100\% éxito), pruebas de
carga hasta 50 usuarios & Cumplido \\
\midrule
Autenticación de usuarios & JWT con bcrypt & Cumplido \\
\midrule
Integridad de los datos & SQLAlchemy como ORM & Cumplido \\
\midrule
Trazabilidad de las operaciones & Logging & Cumplido \\
\midrule
Escalabilidad horizontal & No implementado (SQLite, sin Docker ni balanceo de carga) &
Pendiente \\
\midrule
Alertas por correo electrónico & SMTP solo para recuperación de contraseña, no para alertas
financieras & Parcial \\
\bottomrule
\end{tabular}
\end{table}

El sistema resultante opera como un pipeline secuencial orquestado por FastAPI, donde cada
agente se invoca en orden y los resultados se consolidan en una respuesta única al usuario.
Si bien la arquitectura planificada contemplaba mayor autonomía entre agentes, el enfoque
secuencial resultó suficiente para el prototipo y permitió simplificar la depuración y el
mantenimiento.
